\chapter*{第 7 章\quad 理想流体动力学 课后习题}
\addcontentsline{toc}{chapter}{第 7 章\quad 理想流体动力学 课后习题}
\markboth{第 7 章\quad 理想流体动力学 课后习题}{}

\begin{tishi}
本章共 12 题（教材书 p150--151 的“习 题”段）。本册按教材原顺序排列：
\textbf{第 1 题＝习题 7.1，第 2 题＝习题 7.2，$\cdots$，第 12 题＝习题 7.12}。

本章有一个\textbf{反常识}的结论：考情分析说第七章“概念为主”（E1 slide33），
考点分析说“第六到第八章，少许的填空、判断、选择题会涉及”（E2 slide5），
但翻开真题就会发现——\textbf{本章的计算大题原题密度是全册最高的一章}：
7.3、7.5、7.6、7.9、7.10 五道题都在 2014--2019 年的试题中以计算大题出现过，
7.4 是 2019 年填空题原题，7.1、7.2 的套路是 2016 年、2026 年的计算大题。
所以本章的定位是：\textbf{概念要背，但那几道计算题必须做到“见题就能写”。}

做题要求：给 $\varphi$ 或 $\psi$ 的题，一律按“求偏导定速度 $\to$ 定流线与流量 $\to$ 需要时用伯努利方程求压强”三步写；
判存在性的题，必须把“连续性（管流函数）”与“无旋（管势函数）”两个条件分开写；
叠加类题一律先叠加 $\varphi$、$\psi$（或复势）再求导。
\end{tishi}

\begin{ti}
\sloppy
\emergencystretch=2em

\item 平面不可压缩流体速度分布：$v_x=4x+1$，$v_y=-4y$。
（1）该流动满足连续性方程否？（2）势函数 $\varphi$、流函数 $\psi$ 存在否？（3）求 $\varphi$、$\psi$。
\xstarfull{4}\yiju{E4 2016 年试题计算第 5 题（给速度分布 $u_x=x^2-2x-4y$、$u_y=-2xy+2y$，判连续性、判是否有旋、求势函数与流函数，10 分）；E4 2026 年回忆版计算第 4 题（证明势函数存在性）；E4 2015 年填空第 15 题（速度势函数存在于）、E4 2019 年填空第 3 题、E4 2020 年填空第 4 题、E4 2021 年填空第 8 题、E4 2022 年填空第 7 题共 5 次填空；E5 题库 6-1（无旋流动的定义）、6-2（平面流动具有流函数的条件）、6-4（速度势函数存在的条件）；E2 slide5「第六到第八章……学会套用计算公式即可」}\nandu{中}

\item 平面不可压缩流体速度分布：$v_x=x^2-y^2+x$，$v_y=-(2xy+y)$。
（1）该流动满足连续性方程否？（2）势函数 $\varphi$、流函数 $\psi$ 存在否？（3）求 $\varphi$、$\psi$。
\xstarfull{4}\yiju{E4 2016 年试题计算第 5 题（同一套路，但结论是“有流函数、无势函数”——本题结论相反，正可作对比）；E4 2020 年填空第 4 题（已知 $\varphi=x^2-y^2-x$ 求速度分量）；E4 2021 年填空第 8 题（已知速度势函数求速度分量）；E5 题库 6-2、6-4（流函数与势函数存在的条件）；E2 slide5「学会套用计算公式即可」}\nandu{中}

\item 平面不可压缩流体速度势函数 $\varphi=x^2-y^2-x$，求流场上 $A(-1,-1)$ 及 $B(2,2)$ 点处的速度值及流函数值。
\xstarfull{5}\yiju{E4 2018 年试题计算第 7 题即本题原题（求 $A(-1,-1)$ 及 $B(2,2)$ 点处的速度值及流函数值，10 分）；E4 2020 年填空第 4 题、E4 2021 年填空第 8 题同型（已知势函数或流函数求速度分量）；E6 教材式（7.3）$u_x=\partial\varphi/\partial x$、式（7.8）$u_x=\partial\psi/\partial y$、$u_y=-\partial\psi/\partial x$}\nandu{中}

\item 已知平面流动流函数 $\psi=x+y$，计算其速度、加速度、线变形率 $\varepsilon_{xx}$、$\varepsilon_{yy}$，求出速度势函数 $\varphi$。
\xstarfull{4}\yiju{E4 2019 年试题填空第 3 题即本题原题（已知平面流动流函数 $\psi=x+y$，求 $u_x$、$u_y$、$a_x$、$a_y$ 及速度势函数）；E4 2022 年填空第 7 题（已知 $\psi=-\sqrt{3}x+y$ 求速度分量）；E4 2015 年试题计算第 7 题（同一考点，见第 5 题）；E6 教材式（7.8）、式（3.8a）加速度}\nandu{中}

\item 一平面定常流动的流函数为 $\psi(x,y)=-\sqrt{3}\,x+y$，试求速度分布，写出通过 $A(1,0)$ 和 $B(2,\sqrt{3})$ 两点的流线方程。
\xstarfull{5}\yiju{E4 2015 年试题计算第 7 题即本题原题（求速度分布并写出通过 $A(1,0)$、$B(2,\sqrt{3})$ 的流线方程）；E4 2022 年填空第 7 题即同一流函数（求速度分量）；E6 教材式（7.8）与【例 7.2】——\textbf{本题与教材例 7.2 逐字相同}，属“教材原题直接进场”}\nandu{中}

\item 平面不可压缩流体速度势函数 $\varphi=ax(x^2-3y^2)$（$a<0$），试确定流速及流函数，并求通过连接 $A(0,0)$ 及 $B(1,1)$ 两点的连线的直线段的流体流量。
\xstarfull{5}\yiju{E4 2014 年试题计算第 9 题即本题原题（$\varphi=ax(x^2-3y^2)$，$a<0$，确定流速及流函数，并求通过 $A(0,0)$ 与 $B(1,1)$ 连线直线段的流体流量，8 分）；E6 教材式（7.3）、式（7.6）与流函数特性 3（两流线间单位厚度流量等于流函数值之差）}\nandu{难}

\item 已知复势为：（1）$f(z)=(1+i)z$；（2）$f(z)=(1+i)\ln\dfrac{z+1}{z-4}$；（3）$f(z)=-6iz+\dfrac{24i}{z}$。试分析上述流动的组成。
\xstarfull{3}\yiju{E4 2020 年试题计算第 5 题（一平面均匀流速度大小为 $v_0$、方向与 $x$ 轴成 $\alpha$，求流动的势函数、流函数和复势）；E6 教材 \S7.1.4 复势与复速度（式 7.10、7.11）、\S7.2 基本平面势流（式 7.19、7.28、7.35）、\S7.3 势流叠加；E1 slide33「第七章概念为主」——复势题在真题中出现 1 次，按三星处理}\nandu{难}

\item 已知两个点源布置在 $x$ 轴上相距为 $a$ 的两点，第一个强度为 $2q$ 的点源在原点，第二个强度为 $q$ 的点源位于 $(a,0)$ 处，求流动的速度分布（$q>0$）。
\xstarfull{3}\yiju{E4 2015 年试题计算第 8 题、E4 2019 年试题计算第 5 题（点源与均匀流叠加，见第 10 题）——点源\emph{之间}的叠加真题未直接出现，但“点源速度场 $u_r=\pm q/2\pi r$＋叠加”是必考基本功；E6 教材式（7.20）与 \S7.3.1 叠加原理；E2 slide5「学会套用计算公式即可」}\nandu{中}

\item 如习题 7.9 图所示，平面上有一对等强度为 $\Gamma$（$\Gamma>0$）的点涡，其方向相反，分别位于 $(0,h)$、$(0,-h)$ 两固定点处，同时平面上有一无穷远平行于 $x$ 轴的来流 $v_\infty$，试求合成速度在原点的值。
\tuyi{习题 7.9 图：$xOy$ 平面上，$y$ 轴上、原点上下对称地布置两个点涡，位于 $(0,h)$ 与 $(0,-h)$，两涡强度相等（均为 $\Gamma$）而旋转方向相反；同时在无穷远处有一沿 $+x$ 方向、大小为 $v_\infty$ 的平行来流。所求为原点 $O$ 处的合成速度（因为原点到两涡等距，两个点涡在原点诱导的速度大小都为 $\Gamma/(2\pi h)$，方向均沿 $+x$ 轴）。}
\xstarfull{5}\yiju{E4 2016 年试题计算第 8 题即本题原题（一对等强度反向点涡位于 $(0,h)$、$(0,-h)$，叠加平行来流，求原点合成速度，10 分）；E4 2018 年试题计算第 8 题即本题原题（同图同问，10 分）——两年连考同一道课后习题；E6 教材式（7.30）点涡速度场 $u_\theta=\Gamma/(2\pi r)$、\S7.3.1 叠加原理}\nandu{中}

\item 如习题 7.10 图所示，将速度为 $v_\infty$ 的平行于 $x$ 轴的均匀流和在原点强度为 $q$ 的点源叠加，求叠加后流场中驻点位置及通过驻点的流线方程。
\tuyi{习题 7.10 图：均匀流（速度 $v_\infty$，沿 $+x$）绕过一条“半体”（Rankine 半体）流动，点源置于坐标原点 $O$；图中给出上下两条对称的半体轮廓线、来流箭头，并标注远下游处半体轮廓到 $x$ 轴的距离 $a$ 与点源处的夹角 $\theta$。驻点位于半体前缘、即 $x$ 轴负方向（来流的上游侧）。}
\xstarfull{5}\yiju{E4 2015 年试题计算第 8 题即本题原题（均匀流＋原点强度 $q$ 的点源叠加，求驻点位置及经过驻点的流线方程，12 分）；E4 2019 年试题计算第 5 题即本题原题（同问，10 分）；E4 2019 年试题选择第 3 题（理想均匀流绕流旋转圆柱，分析驻点可能出现的位置）为同一“驻点”考点；E6 教材 \S7.3.1 叠加原理}\nandu{中}

\item 一沿 $x$ 轴正向的均匀流，速度 $U=10\ \mathrm{m/s}$，与一位于原点的点涡相叠加。已知驻点位于点 $(0,-5)$，求：（1）点涡的强度；（2）点 $(0,5)$ 的速度；（3）通过驻点的流线方程。
\xstarfull{3}\yiju{E4 2019 年试题选择第 3 题（理想均匀流绕流旋转圆柱体，分析驻点可能出现的位置）——考的是同一结论“环流把驻点搬到一侧”；本题的具体形式（均匀流＋单点涡）在历年真题中未直接出现；E6 教材式（7.30）--（7.33）点涡的速度场、势函数与流函数}\nandu{难}

\item 一平面势流由点源和点汇叠加而成，点源位于点 $(-1,0)$，其强度 $m_1=20\ \mathrm{m^2/s}$，点汇位于点 $(2,0)$，其强度 $m_2=40\ \mathrm{m^2/s}$，流体密度 $\rho=1.8\ \mathrm{kg/m^3}$。已知流场中点 $(0,0)$ 的压强为 0，试求点 $(0,1)$ 和点 $(1,1)$ 的速度和压强。
\xstarfull{3}\yiju{E4 2020 年试题计算第 5 题（势函数、流函数、复势的构造与速度求解）；E5 题库 6-5（不可压缩流体平面流动中两流线的流函数差值等于单宽体积流量）；E6 教材式（7.20）、式（7.24）、式（7.21）与式（4.9）欧拉积分（无旋流动在全流场适用）；E2 slide5「主要掌握基本知识点概念、计算公式、学会套用计算公式即可」}\nandu{难}

\end{ti}

\begin{banfa}
\textbf{按题号给出的应试提示}\\
\textbf{第 7.1、7.2 题：}两步判据一条也不能省——\\
① \textbf{连续性只“管”流函数}：$\dfrac{\partial v_x}{\partial x}+\dfrac{\partial v_y}{\partial y}=0$（教材式 7.5）成立，只能断定 $\psi$ 存在；
② \textbf{无旋才“管”势函数}：$\dfrac{\partial v_y}{\partial x}=\dfrac{\partial v_x}{\partial y}$（即 $\omega_z=0$，教材式 3.38）成立，才能断定 $\varphi$ 存在。\\
所以课本答案必须写成“满足连续性，$\psi$ 存在；无旋（或有旋），$\varphi$ 存在（或不存在）”，不能只答一个“存在”。
E4 2016 年计算第 5 题的标准答案就是“满足连续性、有流函数，但有旋、不存在势函数”这种“一半有一半没有”的形式，考的就是这句话。\\
\textbf{第 7.1--7.3、7.6 题：}已知 $\varphi$ 求速度＝\textbf{直接求偏导}（$v_x=\dfrac{\partial\varphi}{\partial x}$，$v_y=\dfrac{\partial\varphi}{\partial y}$）；
已知速度求 $\psi$＝用全微分 $d\psi=-v_y\dd x+v_x\dd y$“凑全微分”，最后一定要回代验算
$\dfrac{\partial\psi}{\partial y}=v_x$、$-\dfrac{\partial\psi}{\partial x}=v_y$。流量不要再积分一次：\textbf{两条流线之间的单宽流量＝两端点流函数值之差}。\\
\textbf{第 7.4、7.5 题：}已知 $\psi$ 求速度是“\textbf{换位求偏导}”$v_x=\dfrac{\partial\psi}{\partial y}$、$v_y=-\dfrac{\partial\psi}{\partial x}$，
与已知 $\varphi$ 的情形刚好换位且 $v_y$ 带负号，这是全章第一大错源；
流线方程就是 $\psi(x,y)=C$，把 ``过某点'' 的坐标代进去定 $C$ 即可。\\
\textbf{第 7.7 题：}看到复势先做“\textbf{系数拆解}”：$z$ 的线性项＝均匀流（系数写成 $u_\infty\ee^{-\mathrm{i}\alpha}$ 形式就同时读出速度与方向）；
$\ln(z-z_0)$ 的\textbf{实系数}部分＝点源（正）或点汇（负），\textbf{纯虚系数}部分＝点涡；
$\dfrac{1}{z-z_0}$ 的项＝偶极子；$\ln$、$1/z$ 之外还出现分式 $\dfrac{z+a}{z+b}$ 时，先拆成两个 $\ln$。\\
\textbf{第 7.8--7.12 题：}叠加类题固定四步：\\
① 把各简单流动的 $\varphi$、$\psi$（或复势）\textbf{代数相加}；② 对坐标求偏导得速度场；
③ 驻点令 $v_x=0$、$v_y=0$ 联立求解（先令 $v_y=0$ 往往能立刻定出 $y=0$，再代回求 $x$）；
④ 过驻点的流线：把驻点坐标代入叠加后的 $\psi$ 定出常数 $C$，再写成 $\psi(x,y)=C$。
需要求压强时用 \textbf{无旋流动的伯努利方程（欧拉积分，教材式 4.9）}：$p+\dfrac{1}{2}\rho v^2=$ 常数，全流场通用。
\end{banfa}

\begin{yicuo}
\textbf{本章五个最高频的陷阱：}\\
① \textbf{连续性方程与无旋条件互不代替。}连续性管“$\psi$ 有没有”，无旋管“$\varphi$ 有没有”，
四道真题（2016 计算 5、2019 填空 3、2020 填空 4、2021 填空 8）反复考的就是这一句。
② \textbf{流函数的符号最容易反。}给 $\psi$ 求速度是 $v_x=+\dfrac{\partial\psi}{\partial y}$、$v_y=-\dfrac{\partial\psi}{\partial x}$；
把下面的负号漏掉，速度方向整个反 180$^\circ$。
③ \textbf{势函数与流函数不能并列使用同一组偏导。}$d\varphi=v_x\dd x+v_y\dd y$，而 $d\psi=-v_y\dd x+v_x\dd y$，
两者的符号位置是交叉的——这正是教材式（7.9）的柯西--黎曼条件。\\
④ \textbf{驻点不一定在正半轴。}均匀流＋点源的驻点在 $x=-\dfrac{q}{2\pi v_\infty}$（\textbf{上游}一侧，负 $x$）；
均匀流＋点涡的驻点在被点涡“加速”的那一侧（本题为 $(0,-5)$）。
先写 $v_y=0$ 定方向、再写 $v_x=0$ 定位置，可以避免把驻点找错侧。\\
⑤ \textbf{“两条流线的流量”不是“两点的流量差”。}习题 7.5 中 $A$、$B$ 两点其实落在\emph{同一条}流线上，
两点流函数值之差为 0，所以通过两点连线段的流量为 \textbf{0}——这是教材【例 7.2】埋的坑，也是 2015 年真题的考点。
\end{yicuo}
